\section{Related Work}
\label{sec:related}

\paragraph{In-context medical image segmentation.}
UniverSeg~\cite{butoiUniverSegUniversalMedical2023} popularized ICL
segmentation with a CrossBlock architecture that densely cross-attends
target and context features at full resolution. Iris~\cite{gaoShowSegmentUniversal2025}
instead encodes each context pair into a compact task embedding and
fuses it into the target features through pixel shuffle-unshuffle
operations, avoiding dense cross-attention but performing the
image--label fusion in a single shot with no coarse-to-fine refinement.
Both were designed for 2D slices or small, fixed-size 3D patches; naive
extension to full 3D volumes reintroduces the cubic scaling that
motivated patch- and cascade-based alternatives.

\paragraph{Coarse-to-fine and autoregressive cascades.}
Medverse~\cite{huMedverseUniversalModel2025a} performs next-scale
autoregressive ICL: at each step, the target image and the previous
step's prediction form an ``autoregressive context'' that is processed,
together with the semantic (image, mask) context, by weight-shared U-Net
branches, following the next-scale prediction idea of
VAR~\cite{tian2024visual}. Crucially, every autoregressive step still
performs a full sliding-window pass over the entire volume; only the
\emph{input} to that pass changes with scale, not its spatial extent.
Deep Neural Patchworks~\cite{reisertDeepNeuralPatchworks2022a} similarly
stacks levels of increasing resolution but targets single-model
supervised segmentation rather than in-context, few-shot prediction. In
contrast, PatchICL-3D restricts the sliding window itself to the region
already localized by the previous level, so compute at fine levels
scales with the size of the structure rather than the size of the
volume, and connects levels with register tokens rather than by feeding
back only the rendered prediction.

\paragraph{Image--label fusion mechanisms.}
How a model combines image and label information from the context is a
central design axis. Concatenation-based fusion (Medverse, and earlier
2D CrossBlock-style models) merges the two streams early, before most of
the network's capacity is applied. Iris fuses them once via a
pixel-shuffle/unshuffle operation between context and target. We instead
treat image/label identity as a first-class attention axis, alongside
the spatial-patch axis, allowing the two streams to exchange information
repeatedly across attention layers rather than once at the input or via
a single fusion step. \TODO{position against additional fusion designs
if relevant, e.g. gated cross-attention or FiLM-style conditioning.}

\paragraph{Efficient 3D volumetric backbones.}
Our encoder builds on the nnU-Net residual encoder~\cite{isensee2021nnu}
and 3D axial rotary position embeddings for within-volume and
cross-context self-attention, following the general trend of combining
convolutional feature extraction with attention-based context
aggregation in volumetric segmentation
\cite{hatamizadeh2022unetr,cciccek20163d}. \TODO{expand with any
additional 3D-specific efficiency baselines we compare against.}
